\section{Introduction}\label{sec:intro_background}

Many databases rely on logging to guarantee  consistency and durability 
\cite{DB:OceanBase:VLDB-2022,mysql2013mysql}.
After writing data  to a log file, databases call the costly fsync  to enforce persistence.
Such persistent writes make logging incur severe performance cost~\cite{lee2015waldio}.

Researchers have considered  minimizing the cost of database logging.
At the hardware level, modern solid-state drives (SSDs) are equipped with the power loss protection (PLP). 
PLP enables data supposed to be flushed into persistent media to be
held in the internal cache of SSD~\cite{SSD:PLP:TC-2016,PLP:samsung}.
The hardware time cost  is hence largely reduced for fsync, but 
this  amplifies the relative overhead, i.e., {\em software tax}, imposed by    software layers.
Correspondingly,
to lower the software tax,
researchers
have optimized fsync  at the operating system (OS)
 level~\cite{fsync:rethink:OSDI-2006,fsync:OptFS:SOSP-2013,fsync:FastCommit:ATC-2024,fsync:iJournaling:ATC-2017,fsync:RFLUSH:FAST-2018}.
At the application level, a few databases
  preallocate log files using the fallocate system call (syscall) \cite{DB:OceanBase:VLDB-2022,mysql2013mysql}.
 Preallocation offers two   benefits. First, by reserving contiguous disk space in advance,
 preallocation
 eliminates the   cost 
of runtime block  allocations.
Second, for a preallocated file, invoking fdatasync, rather than fsync, generally does not trigger file system (e.g., Ext4) journaling for metadata.
However, our quantitative study indicates that, even with these joint multi-level optimizations,
logging remains as a serious bottleneck, impacting the performance of database. 
For example, on a low-latency NVMe SSD  with the PLP feature and the latest Ext4,
 disabling all logging actions for MySQL  \cite{mysql2013mysql}
increases its performance by 109.0\%.
Particularly,
the software tax accounts for 60.2\% time for a logging write on average.

On the other hand,
state-of-the-art solutions to mitigate the software tax charged on file I/Os
encounter  adoption barriers due to their implementation complexity and practical limitations. Kernel-bypass techniques, such as SPDK \cite{SSD:SPDK:CloudCom-2017} and NVMeDirect~\cite{SSD:NVMeDirect:HotStorage-2016}, demand extensive reprogramming.
As to approaches involving hardware and/or OS changes, such as Moneta-D \cite{SSD:Moneta-D:ASPLOS-2012} and BypassD~\cite{SSD:BypassD:ASPLOS-2024}, 
their dependencies on   specialized hardware gadgets and intrusive kernel modifications
elevate deployment costs and may introduce potential security vulnerabilities,   limiting their practical utility.

\begin{comment}
 

In this paper, we attempt to exempt the software tax specifically charged on logging I/Os for databases
in an {\em all-in-software} approach.
We attentively examine the structure and usage of a log file.
 A log file, if preallocated and initialized, undergoes virtually no changes in its structural layout.
Generally, neither space reallocation nor alteration of access permissions occurs.
Typically a database keeps reusing the space of its log file in a circular manner.
Consequently, all disk blocks of the file remain stable.
\end{comment}

In this paper, we aim to exempt the runtime software tax charged on
database logging
in an {\em all-in-software} fashion. We are inspired by the observation that a 
 log file, if preallocated, has a stable structure with fixed per-file offset-to-block mappings 
 that can be neatly reorganized and referenced. Moreover,
 writing data to it
  does not alter the file's structural layout. This avoids space  allocation, metadata journaling, 
and access permission changes. 
For non-preallocated log files, we think about if the underlying file system such as Ext4 has implicit 
features that can be used to resemble a preallocation-like effect.

We accordingly develop \odes. \odes builds a shortcut path  that bypasses almost all 
software layers and delivers data directly from database to storage device. 
To this end, it has the following key aspects.
\begin{itemize}[leftmargin=4mm]\setlength{\itemsep}{-\itemsep}
	\item 
	For a preallocated log file,
	\odes extracts and caches the file's offset-to-block mappings 
	in an in-memory structure called the {\em \underline{M}etadat\underline{a} \underline{Co}llector} (Maco). 
	Maco enables \odes to promptly determine target blocks without querying the file's indexing structure (e.g., extent tree in Ext4). 
	\item 
	At runtime, \odes establishes a new shortcut I/O path by leveraging eBPF~\cite{ebpf} 
	and {ioctl}~\cite{ioctlWik84:online}. 
	Referring to per-file mappings kept in Maco, \odes uses eBPF to intercept and redirect 
	persistent writes
	from database to device driver, 
	and pushes down data via ioctl interfaces  to target blocks.
	\item 
	For non-preallocated log files, \odes manages a pool of {\em donor} files to pre-reserve storage spaces.
	Upon an allocation demand to extend a  log file, it exploits Ext4's {\em move-extent} feature \cite{ext4moveextent}
	 to efficiently transfer a range of blocks from a donor file to the log file, which
	is then handled
	in \odes's shortcut I/O path like a preallocated log file. 
\end{itemize}

Ext4 primarily uses move-extent   for  defragmentation \cite{ext4moveextent}.
We    utilize it to online transfer  blocks and 
make a preallocation-like effect for \odes to process persistent writes for non-preallocated log files.
We further optimize \odes as follows.
First, we register preallocated donor files in the Maco  structure
to accelerate  transfers.
Second, we
leverage  io\_uring~\cite{iouring} to do asynchronous transfers outside of the critical path.
Third,
We partition the Maco and donor pool  to be  distributed per CPU core for scalability. 
Last but not the least, we deliver data to storage   by device-specific commands via   ioctl interfaces and enable
compatibility
with  SATA and NVMe SSDs.

\begin{comment}
	we deliver data by
device-specific commands,
such as   {\tt nvme-flush} \cite{nvme-flush-command} for NVMe SSD.
\end{comment}

We  prototype \odes in Linux kernel 6.6.5 without any hardware change. 
Experiments confirm  that 
\odes significantly reduces the   software tax imposed on database logging.
For example,
compared to the stock storage I/O path, \odes achieves 
up to 2.1$\times$ throughput in serving an OLTP write-intensive workload
with  OceanBase~\cite{DB:OceanBase:VLDB-2022} using preallocated log files. 
It is also performant with MySQL whose binary log (binlog) is non-preallocated but gradually growing.
Compared to BypassD lowering the software tax by  
 specialized hardware,
\odes yields much higher performance.

\iffalse 
The rest of this paper is organized as follows. In Section \ref{sec:background}, we present the background knowledge of database logging,
software tax, and eBPF technique. We show our motivational study in Section \ref{sec:motivation}.
We detail the design of \odes in Section \ref{sec:design}. We extensively evaluate \odes against SOTA algorithms
in Section \ref{sec:eval}. We conclude the paper in Section \ref{sec:conclusion}.
\fi

\iffalse
Motivated by this, we develop \odes.
\odes enables a direct I/O path by using eBPF \cite{ebpf} and ioctl \cite{ioctlWik84:online}. In short, after obtaining the  logical block address (LBA), eBPF redirects a logging I/O request from the database to the device driver, in which \odes exploits ioctl interfaces 
 and device-specific commands, such as {\tt nvme-flush} \cite{nvme-flush-command} for NVMe SSD, to deliver data persistently into target disk blocks.

We prototype \odes in Linux kernel 6.6.5.
We chose OceanBase for testing.
For a write-intensive workload, \odes outperforms the status quo I/O path with up to 2.3 $\times$ throughput.
\fi